\chapter{Implementation Details}
\label{chap:implementation-details}

Following the architectural design, the implementation phase translates the requirements into functional code. This chapter details the core implementation strategies of the frontend application and provides an in-depth look at the execution and testing of the primary functionality (F1): User Authentication and Profiles.

\section{Frontend Structure and Routing}
The Angular frontend is structured using Domain-Driven Design (DDD) principles. The application is divided into lazy-loaded feature modules: \texttt{Auth}, \texttt{Feed}, \texttt{Shed}, and \texttt{Events}. Lazy loading ensures that the browser only downloads the code necessary for the current view, significantly improving the initial load time \cite{angularDocs}.

Routing is managed by the Angular Router, which binds specific URL paths to their corresponding components. For example, navigating to \texttt{/login} renders the \texttt{LoginComponent}, while \texttt{/dashboard} renders the \texttt{DashboardComponent}. To enforce security, Route Guards (\texttt{AuthGuard}) are implemented to prevent unauthenticated users from accessing protected routes, redirecting them back to the login page if they lack a valid session token.

\section{Realizing Functionality F1: User Authentication}
User Authentication (F1) is the foundational functionality of the NEST platform, ensuring that community interactions remain secure and attributable to verified residents.

\subsection{Implementation of the Login Flow}
The login process relies on Reactive Forms in Angular to capture user input safely and validate it synchronously (e.g., ensuring the email format is correct and the password meets length requirements) before submission.

When a user submits the form, the \texttt{AuthService} dispatches an NgRx action containing the credentials. An NgRx Effect intercepts this action, making an HTTP POST request to the Express backend (\texttt{/api/auth/login}). 
If the backend validates the credentials against the hashed password in the database, it returns a JSON Web Token (JWT). The Effect then dispatches a success action, updating the Global Store with the user's profile and token, and the Router navigates the user to the main feed.

\subsection{HTTP Interceptors}
To maintain the authenticated state across the application, an \texttt{AuthInterceptor} is implemented. This service intercepts all outgoing HTTP requests from the frontend and automatically attaches the JWT to the \texttt{Authorization} header. This eliminates the need to manually append the token to every single API call, reducing code duplication and preventing security oversights.

\section{Testing Functionality F1}
Validating the authentication flow is critical. The testing strategy involves both automated unit testing and manual execution verification.

\subsection{Execution and Verification}
To demonstrate F1 as an executable functionality, the following manual test case is performed:
1. \textbf{Setup:} The backend server and database are running locally. The Angular development server is launched.
2. \textbf{Action:} The user navigates to \texttt{localhost:4200/login}, enters valid credentials, and clicks 'Log In'.
3. \textbf{Expected Result:} The application briefly displays a loading spinner, the JWT is stored in the browser's \texttt{localStorage}, and the view transitions to the Community Feed. A welcome toast notification is displayed.
4. \textbf{Negative Test:} The user enters an incorrect password. The system must prevent navigation and display an inline error message ("Invalid credentials") without crashing.

\textbf{REPLACE HERE WITH: [A screenshot of the application in execution during the F1 phase. Ideally, show the login screen with validation errors visible, or the successful transition to the user profile/dashboard, demonstrating the executable state of F1.]}

\section{Implementation of Secondary Modules}
While Authentication forms the base, the subsequent modules follow a similar implementation pattern:
\begin{itemize}
    \item \textbf{Community Feed:} Utilizes Angular's \texttt{AsyncPipe} to subscribe directly to the NgRx Store, rendering posts dynamically as they arrive without manual subscription management.
    \item \textbf{Shared Shed:} Implements complex reactive forms for item creation, including image upload handling, and communicates with the \texttt{ShedApiService} to sync state with the backend database.
\end{itemize}
